\newpage
\chapter{Выводы по работе}

\section{Выводы по исследованию амплитудной модуляции}

Для контрольного режима с $m = 0.8$ экспериментальные значения спектральных составляющих и длины вектора (таблица~\ref{tab:AM_measurements}) полностью совпали с расчетами из раздела~\ref{sec:AM_calculation}.

Фазовый портрет АМ-сигнала представляет собой прямую горизонтальную линию вдоль вещественной оси $I$.
Это объясняется тем, что амплитуда сигнала изменяется во времени, а его фаза остается постоянной.

Длина отрезка на IQ-плоскости определяется удвоенным размахом изменения амплитуды огибающей:
\begin{equation*}
    L = U_{\max} - U_{\min} = U_0 (1 + m) - U_0 (1 - m) = 2 \cdot m \cdot U_0.
\end{equation*}

Поскольку $U_0 = 1.0$ В, то длина линии численно равна $L = 2 \cdot m$ и возрастает линейно от $0.6$ В при $m = 0.3$ до $2.0$ В при $m = 1.0$.

Расчет по формуле~\eqref{eq:AM_modulation_index} для коэффициента модуляции $m$ полностью совпал с формой огибающей сигнала.

\section{Выводы по исследованию частотной модуляции}

При $m = 0.2$ измеренные значения составили: несущая $0.99$ В (теоретически $1.0$ В), первые боковые $0.10$ В (теоретически $0.10$ В), ширина полосы $2.0$ кГц (теоретически $2.0$ кГц).

При $m = 5.0$ измеренные значения составили: несущая $0.18$ В (теоретически по Бесселю $|J_0(5)| = 0.178$ В), первые боковые $0.33$ В (теоретически $|J_1(5)| = 0.328$ В), вторые боковые $0.05$ В (теоретически $|J_2(5)| = 0.047$ В).
Расхождение между экспериментальными и теоретическими значениями объясняется округлением и не превышает половины процента.

При увеличении индекса модуляции $m$ от $0.2$ до $5.0$ ширина полосы частот возрастает строго монотонно от $2.0$ кГц до $14.0$ кГц.
В отличие от АМ-сигнала, при ЧМ происходит постоянное перераспределение энергии между несущей и боковыми составляющими.
В частности, при $m = 2.6$ амплитуда несущей практически подавляется до $0.10$ В.

Амплитуда огибающей ЧМ-сигнала остается постоянной, поэтому траектория вектора всегда лежит на окружности с радиусом $1$ В.
Угловой сектор развертки вектора равен $\Delta\theta = 2m$ радиан.
Дуга полностью замыкается в кольцо при $2m \geqslant 2\pi$, что соответствует $m \geqslant \pi \approx 3.14$.
В эксперименте замыкание кольца зафиксировано на шаге $m = 3.4$.

\section{Выводы по исследованию фазовой модуляции}

Визуально фазовая модуляция очень похожа на частотную, так как огибающие обеих модуляций постоянны по амплитуде, а плотность периодов высокочастотных колебаний сжимается и растягивается.
Однако, отличие заключается в фазовом сдвиге: при гармоническом сообщении максимум девиации частоты ФМ-сигнала сдвинут на четверть периода относительно ЧМ.
Это связано с тем, что мгновенная частота ФМ-сигнала пропорциональна производной от фазы $\omega(t) = \omega_0 + m\varOmega\cos{(\varOmega t)}$.

Фазовый портрет ФМ-сигнала представляет собой дугу окружности единичного радиуса.
Амплитуда вектора не меняется, изменяется только его фазовый угол.
Угловой размах сектора составляет $\Delta\theta = 2m = 4.0$ рад ($\approx 229^\circ$), поэтому дуга не замыкается при $m = 2.0$.

При увеличении индекса модуляции до $m = 4.0$ угловой размах сектора составляет $\Delta\theta = 8.0$ рад ($\approx 458^\circ > 360^\circ$).
Это превышает полный оборот окружности, поэтому дуга на комплексной плоскости $IQ$ полностью замыкается в кольцо.

\section{Выводы по исследованию двухпозиционной фазовой манипуляции}

В моменты смены бита на высокочастотной осциллограмме наблюдается мгновенный скачкообразный переворот фазы сигнала на $180^\circ$.
Длительность одного символа составляет $T_s = 1$ мс, что соответствует символьной скорости манипуляции:
\begin{equation*}
    R_s = \frac{1}{T_s} = \frac{1}{1 \cdot 10^{-3} \text{ с}}  = 1000\text{ Бод} = 1.0\text{ кБод}.
\end{equation*}

Поскольку каждый символ кодирует ровно один бит информации, то информационная скорость численно равна символьной скорости: $R = R_s = 1.0$ кбит/с.

Созвездие ФМ-2 состоит из двух точек, симметрично расположенных на синфазной оси $I$ в координатах $(-1.0\text{ В}; 0)$ для логического «0» и $(+1.0\text{ В}; 0)$ для логической «1».
Квадратурная составляющая $Q$ равна нулю.
Траектория перехода между точками созвездия представляет собой прямую линию, проходящую через центр координат $(0;0)$.

Спектр является непрерывным.
Центральный пик спектра центрирован на частоте несущей $f_0 = 12.0$ кГц, а его ширина составляет ровно удвоенную тактовую частоту:
\begin{equation*}
    \Delta f = 2 \cdot R_s = 2 \cdot 1.0\text{ кГц} = 2.0\text{ кГц}.
\end{equation*}

Нули спектра приходятся на частоты $11$ кГц и $13$ кГц.

\section{Выводы по исследованию четырехпозиционной фазовой манипуляции}

Форма и ширина спектра ФМ-4 практически совпадают со спектром ФМ-2: центральный пик центрирован на частоте несущей $f_0 = 12.0$ кГц, а его ширина составляет $\Delta f = 2.0$ кГц.
Это связано с тем, что символьная скорость манипуляции осталась прежней: $R_s = 1.0$ кБод, длительность символа $T_s = 1$ мс, следовательно, занимаемая полоса частот не изменилась.

Созвездие ФМ-4 содержит четыре точки, симметрично расположенные на комплексной плоскости $IQ$ в каждом из квадрантов.

Углы фазового наклона векторов составляют $45^\circ$, $135^\circ$, $225^\circ$ и $315^\circ$.

Минимальный фазовый сдвиг между соседними точками составляет $90^\circ$ ($\pi/2$ рад).

Показания измерительного табло зафиксировали скачок информационной скорости с $1.0$ кбит/с для ФМ-2 до $2.0$ кбит/с для ФМ-4.
Скорость удвоилась, так как каждый символ кодирует не один бит, а два бита информации.
Это доказывает двукратный рост эффективности ФМ-4 по сравнению с ФМ-2 при сохранении той же полосы пропускания канала.

\section{Выводы по исследованию двухпозиционной частотной манипуляции}

На экране анализатора спектра зафиксированы два выраженных пика, расположенных на частотах $f_1 = 9.0$ кГц и $f_2 = 15.0$ кГц.
Наличие множества мелких промежуточных пиков объясняется структурой спектра и взаимным наложением боковых составляющих.

Частотное расстояние между двумя главными пиками составляет:
\begin{equation*}
    \Delta F = |f_2 - f_1| = |15.0 - 9.0| = 6.0\text{ кГц}.
\end{equation*}

При этом девиация частоты относительно центральной несущей $f_0 = 12.0$ кГц равна:
\begin{equation*}
    \Delta f = \frac{\Delta F}{2} = \frac{6.0}{2} = 3.0\text{ кГц}.
\end{equation*}

Разнос частот существенно превышает тактовую символьную скорость передачи $R_s = 1.0$ кБод.
Такой широкий разнос обеспечивает надежную устойчивость к помехам и шумам, однако требует значительной полосы пропускания радиочастотного спектра.

\section{Выводы по исследованию четырехпозиционной частотной манипуляции}

В отличие от ЧМ-2, где наблюдалось 2 пика, в спектре ЧМ-4 зафиксировано 4 выраженных пика, соответствующих четырем дибитам.
Пики расположены симметрично относительно несущей частоты $f_0 = 12.0$ кГц на частотах $f_1 = 7.5$ кГц, $f_2 = 10.5$ кГц, $f_3 = 13.5$ кГц и $f_4 = 16.5$ кГц.

Частотный шаг между соседними пиками составляет:
\begin{equation*}
    \delta f = |f_2 - f_1| = |f_3 - f_2| = |f_4 - f_3| = 3.0\text{ кГц}.
\end{equation*}

При переходе от ЧМ-2 к ЧМ-4 информационная скорость передачи возросла в 2 раза с $1.0$ кбит/с до $2.0$ кбит/с за счет кодирования двух битов информации в одном символе.
Однако для передачи четырех состояний потребовалось расширить полосу частот между крайними пиками с $6.0$ кГц до $9.0$ кГц.
Получается, что удельная спектральная эффективность возрастает незначительно с $0.17$ до $0.22$ бит/(с$\cdot$Гц).
Это подтверждает теоретическое свойство частотной манипуляции: увеличение скорости всегда требует расширения полосы радиочастот.
А в фазовой модуляции скорость удваивается без расширения спектра.

\section{Выводы по исследованию амплитудно-фазовой манипуляции (АФМ-16)}

В сигналах чисто фазовой манипуляции, таких как ФМ-2 и ФМ-4, огибающая была строго постоянной, а модуляция осуществлялась только скачками фазы.
В сигнале АФМ-16 управление несущей происходит одновременно по двум параметрам --- фазе и амплитуде.
На осциллограмме (рис.~\ref{fig:APSK-16}) четко наблюдаются скачки между двумя дискретными уровнями амплитуды, соответствующими радиусам внутреннего и внешнего колец.

Всего у созвездия 16 точек.
Все они распределены по двум окружностям:
\begin{itemize}
    \item Внутреннее кольцо содержит ровно 4 точки, симметрично лежащие в четырех квадрантах с шагом фазы $90^\circ$.
    Точки расположены под углами $45^\circ$, $135^\circ$, $225^\circ$, $315^\circ$;
    \item Внешнее кольцо содержит 12 точек, равномерно распределенных по окружности с шагом фазы $30^\circ$ или $\frac{\pi}{6}$ радиан.
\end{itemize}

Все точки на каждом из колец имеют одинаковую амплитуду.
Такое кольцевое распределение делает сигнал АФМ-16 более устойчивым к нелинейным искажениям, например, усилителей мощности спутниковых передатчиков.

Каждый радиосимвол кодирует 4 бита информации ($M = 2^4 = 16$), следовательно, информационная скорость выросла до $4.0$ кбит/с.
Это в 4 раза выше по сравнению с ФМ-2 и ФМ-4 при сохранении той же полосы частот радиосигнала в эфире.

\section{Выводы по исследованию амплитудно-фазовой манипуляции (АФМ-32)}

Внешний вид и ширина главного пика спектра АФМ-32 практически не отличаются от спектра АФМ-16.
Это связано с тем, что символьная скорость системы осталась строго фиксированной: $R_s = 1.0$ кБод, длительность символа $T_s = 1.0$ мс.
Следовательно, при переходе от АФМ-16 к АФМ-32 рост пропускной способности достигается без расширения занимаемой полосы радиочастот.

Всего у созвездия 32 точки.
Все они распределены по трем кольцам:
\begin{itemize}
    \item Внутреннее кольцо содержит 4 точки, по одной в квадранте, с шагом $90^\circ$;
    \item Среднее кольцо содержит 12 точек с равномерным шагом $30^\circ$;
    \item Внешнее кольцо содержит 16 точек с равномерным шагом $22.5^\circ$ или $\frac{\pi}{8}$ радиан.
    На осциллограмме (рис.~\ref{fig:APSK-32}) этой геометрии соответствуют уже три дискретных уровня огибающей радиосигнала.
\end{itemize}

Зафиксирована информационная скорость $R$, равная $5.0$ кбит/с.
Однако, как мы знаем, символьная скорость системы составляет $R_s = 1.0$ кБод.
Один символ переносит 5 бит данных, $n = \log_2(32)$.
Тогда информационная скорость рассчитывается так:
\begin{equation*}
    R = R_s \cdot n = 1.0 \text{ кБод} \cdot 5 \text{ бит/символ} = 5 \text{ кбит/с}.
\end{equation*}

Это подтверждает пятикратное повышение спектральной эффективности по сравнению с базовой манипуляцией ФМ-2.

\section{Выводы по исследованию квадратурной амплитудной манипуляции (КАМ-16)}

В отличие от АФМ-16, где точки группировались по двум кольцам, в КАМ-16 созвездие представляет собой симметричную квадратную матрицу $4 \times 4$, содержащую ровно 16 сигнальных точек.

В АФМ-16 все точки одного кольца имели одинаковую амплитуду.
В КАМ-16 4 внутренние точки имеют наименьшую амплитуду, 8 промежуточных боковых --- среднюю, а 4 крайние угловые --- максимальную.

Форма и ширина спектра КАМ-16 выглядят похожими на спектр АФМ-16.
В обоих методах символьная скорость одинакова: $R_s = 1.0$ кБод, $T_s = 1.0$ мс.

По данным табло Д, информационная скорость составила $4.0$ кбит/с, что ровно в 4 раза выше, чем у ФМ-2.
Это обусловлено тем, что в КАМ-16 каждый передаваемый символ кодирует 4 бита информации, $n = \log_2(16) = 4$.
Символьная скорость манипуляции не изменилась.

Также на табло зафиксировано значение пиковой амплитуды $U_m \approx 1.41$ В.
Значение в точности совпало с теоретическим расчетом из раздела~\ref{sec:qam-16--qam-32_measurements}:
\begin{equation*}
    U_m = \sqrt{I_{\max}^2 + Q_{\max}^2} = \sqrt{1.0^2 + 1.0^2} = \sqrt{2} \approx 1.414 \text{ В.}
\end{equation*}

\section{Выводы по исследованию квадратурной амплитудной манипуляции (КАМ-32)}

В КАМ-16 созвездие представляло собой стандартный квадрат $4 \times 4$, 16 точек.
В КАМ-32 для размещения 32 точек нужна расширенная сетка $6 \times 6$, 36 точек.
Однако построение квадрата привело бы к избыточности и резкому росту пиковой мощности на 4 крайних углах.
Поэтому в схеме реализовано созвездие-крест: 4 крайние угловые точки удалены, $36 - 4 = 32$.
Это делает созвездие более компактным и существенно повышает его помехоустойчивость.

Спектр КАМ-32 по форме и занимаемой полосе частот практически не отличается от АФМ-32.
Причина в том, что в обеих схемах тактовая символьная скорость манипуляции одинакова.
Разница состоит в способе кодирования плоскости --- кольцо и сетка.

Зафиксированный на табло Д максимальный уровень радиосигнала составил $U_m \approx 1.30$ В.
Это значение совпадает с теоретическим расчетом из раздела~\ref{sec:qam-16--qam-32_measurements}:
\begin{equation*}
    U_m = \sqrt{I_{\max}^2 + Q_{\max}^2} = \sqrt{(\frac{5}{4.5})^2 + (\frac{3}{4.5})^2} = \frac{\sqrt{34}}{4.5} \approx 1.296 \text{ В.}
\end{equation*}

Удаление 4 углов позволило удержать пиковое напряжение на уровне $1.30$ В вместо $1.57$ В, который бы возник в полном квадрате $6 \times 6$.